\section*{الفصل \ref{ch:b1:flowering-plant-organization} --- التنظيم الوظيفي عند النبات الزهري}

\begin{solution}{exo:b1:flowering-plant-organization:1}
\omterm{def:b1:flowering-plant-organization:organs}{الجذر}: التثبيت، وامتصاص الماء والشوارد المعدنية. الساق: حمل \omterm{def:b1:flowering-plant-organization:organs}{الأوراق}
نحو الضوء، والنقل بين \omterm{def:b1:flowering-plant-organization:organs}{الجذور} \omterm{def:b1:flowering-plant-organization:organs}{والأوراق}. \omterm{def:b1:flowering-plant-organization:organs}{الورقة}: التركيب الضوئي
(والتبادل الغازي، والنتح).
\end{solution}

\begin{solution}{exo:b1:flowering-plant-organization:2}
الوحدة \omterm{def:b1:flowering-plant-organization:organs}{عقدة} واحدة بورقتها وبرعمها الإبطي، مع \omterm{def:b1:flowering-plant-organization:organs}{السلامية} التي تحتها.
\omterm{prop:b1:flowering-plant-organization:modular}{والنمو غير المحدود}: لا قياس بالغ؛ فالبراعم القمية والإبطية تظل تضيف
الوحدات ما دام النبات حيًا.
\end{solution}

\begin{solution}{exo:b1:flowering-plant-organization:3}
البشرة: حية؛ تغطية، \omterm{def:b1:flowering-plant-organization:tissues}{وأديمة}، وثغور، وأشعار جذرية. \omterm{def:b1:flowering-plant-organization:tissues}{البرنشيم}: حي؛
تركيب ضوئي وادخار. \omterm{def:b1:flowering-plant-organization:tissues}{الكولنشيم}: حي؛ دعامة مرنة. \omterm{def:b1:flowering-plant-organization:tissues}{السكلرنشيم}: ميت؛
دعامة صلبة. خلايا \omterm{def:b1:flowering-plant-organization:tissues}{الخشب} الناقلة: ميتة؛ نقل الماء. \omterm{def:b1:flowering-plant-organization:tissues}{الأنابيب الغربالية}
في \omterm{def:b1:flowering-plant-organization:tissues}{اللحاء}: حية (مع خلاياها المرافقة)؛ نقل السكر.
\end{solution}

\begin{solution}{exo:b1:flowering-plant-organization:4}
$e^{-0.9} = 0.41$، أي \qty{41}{\%}؛ و$e^{-3} = 0.05$، أي
\qty{5}{\%}.
\end{solution}

\begin{solution}{exo:b1:flowering-plant-organization:5}
\omterm{def:b1:flowering-plant-organization:organs}{جذر} (أسطوانة مركزية واحدة، \omterm{prop:b1:flowering-plant-organization:stemroot}{وبشرة داخلية}، \omterm{def:b1:flowering-plant-organization:tissues}{وخشب} \omterm{def:b1:flowering-plant-organization:tissues}{ولحاء} متناوبان على
نصف قطر) لثنائي الفلقة (أذرع \omterm{def:b1:flowering-plant-organization:tissues}{خشب} قليلة، ولا لب مركزي).
\end{solution}

\begin{solution}{exo:b1:flowering-plant-organization:6}
العارضة المنحنية أشد إجهادًا عند سطحها ولا إجهاد على محورها، فالمادة
الموضوعة في حلقة محيطية تقاوم الانحناء بأقل كتلة (أي أنبوب). أما
الكبل المشدود فإجهاده منتظم على مقطعه وهو ينحني بحرية؛ فيحمل خصل
مركزي الشدّ بينما تبقى القشرة مرنة، وتستطيع \omterm{def:b1:flowering-plant-organization:organs}{الجذور} الجانبية أن تخرج
عبرها.
\end{solution}

\begin{solution}{exo:b1:flowering-plant-organization:7}
$1 - e^{-2.4} = 0.91$؛ وعند $L = 6$ يكون $1 - e^{-3.6} = 0.97$:
فطبقتان إضافيتان من الورق تضيفان \qty{6}{\%} من الضوء وتكلفان
\qty{50}{\%} ورقًا زائدًا.
\end{solution}

\begin{solution}{exo:b1:flowering-plant-organization:8}
$S = \pi d L$: فيكون
$d = 240/(\pi\times 623\,000) = \qty{1.2e-4}{m}$، أي نحو
\qty{0.12}{mm}. و$(240 + 400)/5 = 128$.
\end{solution}

\begin{solution}{exo:b1:flowering-plant-organization:9}
سطح \omterm{def:b1:flowering-plant-organization:organs}{الجذر} لكل مليمتر: $\pi\times 0.4 = \qty{1.26}{mm^2}$. والأشعار
لكل مليمتر: $250\times\pi\times 0.010\times 0.8 = \qty{6.3}{mm^2}$.
والمجموع $7.5/1.26 = 6$: أي إن الأشعار تضاعف السطح ست مرات.
\end{solution}

\begin{solution}{exo:b1:flowering-plant-organization:10}
في الغابة تكون البراعم الإبطية السفلى مظللة؛ فالوحدات التي كانت
ستصنعها لا تسدد ثمن نفسها، فتبقى ساكنة أو تموت فروعها، ويُخصص النمو
للبرعم القمي المتسابق صعودًا نحو الضوء. أما في الحقل فكل برعم مضاء
وكل فرع يسدد، فيمتلئ التاج منخفضًا. والكلفة: جذع طويل نحيل استثمار
كبير في \omterm{def:b1:flowering-plant-organization:tissues}{خشب} غير منتج، معرض للريح؛ والكسب: أن يوافق النبات شكله مع
موضع الضوء، دون أي خطة مسبقة.
\end{solution}

\begin{solution}{exo:b1:flowering-plant-organization:11}
الحفرة والأشعار تحبس طبقة من الهواء الساكن الرطب فوق الفتحة؛ فينتشر
تدرج البخار من الحيزات الهوائية إلى الهواء المتحرك على مسار أطول
فيهبط النتح. وعلى ثاني أكسيد الكربون أن ينتشر داخلًا على المسار
الأطول نفسه، فيهبط الأخذ أيضًا، وإن كان هبوطًا أقل، لأن تدرج
$\mathrm{CO_2}$ يحدده الجو واستهلاك \omterm{def:b1:flowering-plant-organization:organs}{الورقة}. وفي الصحراء يكون الماء
هو المورد المحدد فتنقذ المقايضة حياة النبات؛ وفي غابة مطيرة يكون
الماء مجانيًا ويكون كسب الكربون هو ما تكافئ عليه المنافسة.
\end{solution}

\begin{solution}{exo:b1:flowering-plant-organization:12}
ليس للنبات \omterm{def:b1:mammal-organization:compartments}{سائل خارج خلوي} منظَّم يُحفظ عند حرارة وتركيب ثابتين بين
خلاياه والخارج: فخلاياه تلاقي ماء التربة والهواء وضوء الشمس مباشرة،
وحرارته حرارة الهواء. ولذلك تنظم كل خلية محتوياتها (شوارد، وماء،
وpH) عبر غشائها وجدارها؛ وينظم النبات تبادلاته عند السطح (الثغور
تنفتح وتنغلق، \omterm{prop:b1:flowering-plant-organization:stemroot}{والبشرة الداخلية} تنتقي ما يدخل \omterm{def:b1:flowering-plant-organization:tissues}{الخشب}) ويكيّف جسمه
بالنمو (\omterm{def:b1:flowering-plant-organization:organs}{جذور} أكثر في التربة الجافة، \omterm{def:b1:flowering-plant-organization:organs}{وأوراق} أكثر في الظل). فهو ينظم،
لكن التنظيم عند السطوح وفي الشكل، لا في سائل خاص.
\end{solution}

\begin{solution}{pb:b1:flowering-plant-organization:1}
\textbf{1.} $5\times 100 = \qty{500}{m^2}$.
\textbf{2.} $e^{-2.5} = 0.082$ يبلغ الأرض؛ و\qty{92}{\%} معترَض.
\textbf{3.} عند $L = 2$: $1 - e^{-1} = \qty{63}{\%}$؛ وعند $L = 8$:
$1 -
e^{-4} = \qty{98}{\%}$.
\textbf{4.} من 2 إلى 5 يكسب التاج \qty{29}{\%} من الضوء مقابل
\qty{300}{m^2} من الورق؛ ومن 5 إلى 8 سيكسب \qty{6}{\%} مقابل
\qty{300}{m^2} أخرى تكلف بناءً وإمدادًا بالماء بالقدر نفسه: \omterm{def:b1:flowering-plant-organization:organs}{فالأوراق}
الزائدة لا تسدد.
\textbf{5.} $500/0.005 = \num{100000}$ \omterm{def:b1:flowering-plant-organization:organs}{ورقة}.
\textbf{6.} $500\times 10^6\,\mathrm{mm^2}\times 150 = \num{7.5e10}$
\omterm{def:b1:flowering-plant-organization:tissues}{ثغر}.
\textbf{7.} $\qty{50}{\micro m^2}$؛ ولكل mm$^2$:
$150\times 50\times
10^{-6} = \qty{7.5e-3}{mm^2}$، أي \qty{0.75}{\%}
من السطح.
\textbf{8.} $500\times 4\times 10^{-6} = \qty{2e-3}{mol/s}$، أي
\qty{88}{mg/s}.
\textbf{9.} $0.088\times 43\,200 = \qty{3.8}{kg}$ من
$\mathrm{CO_2}$، تحوي $3.8\times 12/44 = \qty{1.04}{kg}$ من الكربون.
\textbf{10.} $150\times 1.04 = \qty{156}{kg}$ من الكربون؛ ونصفه، أي
\qty{78}{kg}، إلى \omterm{def:b1:flowering-plant-organization:tissues}{خشب} فيه \qty{50}{\%} كربون: أي \qty{156}{kg} من
\omterm{def:b1:flowering-plant-organization:tissues}{الخشب}.
\textbf{11.} $500\times 10^{-3} = \qty{0.5}{mol/s}$، أي \qty{9}{g/s}.
\textbf{12.} $9\times 43\,200 = \qty{389}{kg}$، أي نحو \qty{390}{L}.
\textbf{13.} $0.5/(2\times 10^{-3}) = 250$ جزيء ماء لكل
$\mathrm{CO_2}$. فالفتحة التي تُدخل $\mathrm{CO_2}$ تُخرج الماء؛
والهواء الداخلي مشبع بينما $\mathrm{CO_2}$ ليس إلا \qty{0.04}{\%} من
الهواء الخارجي، فيكون تدرج البخار إلى الخارج أشد انحدارًا بكثير من
تدرج $\mathrm{CO_2}$ إلى الداخل.
\textbf{14.} $\qty{2}{mm}\times\qty{100}{m^2} = \qty{200}{L}$: أي
نصف يوم.
\textbf{15.} \qty{100}{m^3}؛ و$5\times 100 = \qty{500}{km}$.
\textbf{16.} $\pi\times 0.5\times 10^{-3}\times 5\times 10^5 =
\qty{785}{m^2}$.
\textbf{17.} $200\times 5\times 10^8\,\mathrm{mm} = \num{1e11}$ شعرة.
\textbf{18.} كل واحدة $\pi\times 10^{-5}\times 5\times 10^{-4} =
\qty{1.57e-8}{m^2}$؛ والمجموع \qty{1570}{m^2}.
\textbf{19.} $785 + 1570 = \qty{2355}{m^2}$، أي $4.7$ أمثال مساحة
\omterm{def:b1:flowering-plant-organization:organs}{الأوراق}.
\textbf{20.} الحجم الواقع ضمن \qty{1}{mm}:
$\pi\times(10^{-3})^2\times 5
\times 10^5 = \qty{1.6}{m^3}$، أي
\qty{1.6}{\%} من التربة.
\textbf{21.} الشيلم: $\pi\times 10^{-6}\times 6.2\times 10^5 =
\qty{1.95}{m^3}$ من صندوق سعته \qty{0.05}{m^3}، أي \qty{3900}{\%}:
فكل نقطة من التربة على بعد مليمتر من أربعين جذرًا. فالشيلم يستكشف
تربته استكشافًا أدق بأكثر من ألفي مرة --- نبتة محصول بميزانية أربعة
أشهر مقابل شجرة تستكشف مئة متر مكعب.
\textbf{22.} $1000 + 2355 = \qty{3355}{m^2}$؛ أي $\qty{34}{m^2}$
لكل متر مربع من الأرض.
\textbf{23.} عند الإنسان: $130/2 = 65$؛ فنسبة البلوطة نصف نسبة
السطوح المطوية عند الإنسان، دون أي طي.
\textbf{24.} البلوطة تضيف وحدات وأوراقًا وجذورًا كل سنة، فتنمو
سطوحها مع عمرها؛ أما سطوح الثديي \omterm{def:b1:mammal-organization:organ}{فأعضاء} ثابتة، تكتمل عند القياس
البالغ، والسبيل الوحيد إلى رفع نسبتها كان طيها في أثناء النمو.
\textbf{25.} نحو \qty{34}{m^2} من \omterm{prop:b1:organism-environment:surfaces}{سطح التبادل} لكل متر مربع من
الأرض: \qty{10}{m^2} من الورق (بوجهيه) فوقها و\qty{24}{m^2} من
\omterm{def:b1:flowering-plant-organization:organs}{الجذور} والأشعار الجذرية تحتها.
\end{solution}
